\section{Quantitative Performance Evaluation}
\label{sec:quantitative_performance}

This section separates the previously reported benchmark from a new controlled
performance matrix for the restructured experiments.  The legacy measurements
are retained because they document the computational behavior of the original
case studies.  The new measurements use a uniform protocol and are reported in
separate tables so that concurrent production-batch timings are not presented as
controlled benchmark observations.

\subsection{Legacy Benchmarks}

The numerical values in Table~\ref{tab:quantitative-performance-legacy} are
retained unchanged from the original submission.  Each configuration was run
five times in the 94-region Porto Alegre model, and the table reports the maximum
blob count, simulated cycles, mean runtime, runtime per cycle, and two historical
memory measurements.  MemProfile measured the peak resident memory of the
Python process, whereas TraceMalloc measured traced Python allocations.  The
references in the case-study column have been updated to the restructured
results in Sections~\ref{sec:results_routine_mobility} and
\ref{sec:results_large_scale_events_epidemics}.

The random-walk configuration with $levy\_scale=2$ is the legacy baseline.  It
completed its single simulated cycle in 5.211~s on average.  Movement
restrictions in Figure~\ref{fig:routine_levy_restriction} reduced the runtime per
cycle to 4.795~s at a restriction of 0.75 while leaving peak memory nearly
unchanged.  Large-scale gatherings in
Figure~\ref{fig:large_scale_event_displacements_94} increased runtime as more
gather actions drew population from multiple regions.  The epidemic and
vaccination experiments in Figures~\ref{fig:epidemic_r0_94},
\ref{fig:vaccination_results_94}, and
\ref{fig:vaccination_restriction_results_94} were more demanding because added
traceable states produced additional population splitting and fewer merge
opportunities.  The largest legacy blob count, 23,355, occurred in VaccC and
coincided with the largest runtime per cycle, 33.455~s.

All legacy and new benchmarks describe the current machine: an AMD Ryzen 5
5600X CPU, an NVIDIA GeForce RTX 2060 GPU, 16~GB DDR4-2666 CL19 memory, and a
KINGSTON SNV3S1000G SSD, using Python 3.13.6 on Microsoft Windows 11 Education
(version 10.0.26200).  The corrected Windows version is reported here; the
original text contained an additional zero.

\begin{table}[htbp]
\footnotesize
\centering
\caption{Retained legacy quantitative performance evaluation.  Numerical values
are unchanged from the original submission.  Each row reports the maximum blob
count, cycles, mean runtime over five runs, runtime per cycle, peak process RSS,
and peak traced Python allocation.}
\label{tab:quantitative-performance-legacy}
\resizebox{\textwidth}{!}{%
\begin{tabular}{@{}clrrrrrr@{}}
\toprule
Case study & Configuration & Max blobs & Cycles & Mean runtime (s) & s/cycle & MemProfile (MiB) & TraceMalloc (MiB) \\
\midrule
Baseline random walk (Section~\ref{sec:results_routine_mobility})
  & $levy\_scale=2$ & 2857 & 1 & 5.211 & 5.211 & 234.912 & 58.565 \\
\midrule
\multirow{3}{*}{Movement restriction (Figure~\ref{fig:routine_levy_restriction})}
  & Restriction $=0.25$ & 2871 & 1 & 5.074 & 5.074 & 234.957 & 58.577 \\
  & Restriction $=0.50$ & 2864 & 1 & 4.958 & 4.958 & 234.935 & 58.527 \\
  & Restriction $=0.75$ & 2883 & 1 & 4.795 & 4.795 & 235.234 & 58.516 \\
\midrule
\multirow{5}{*}{Large-scale events (Table~\ref{tbl:large_scale_event_displacements_94})}
  & Average of 1 event & 2671 & 1 & 6.144 & 6.144 & 235.355 & 58.767 \\
  & Average of 2 events & 2651 & 1 & 7.986 & 7.986 & 235.297 & 58.892 \\
  & Average of 3 events & 2651 & 1 & 9.910 & 9.910 & 235.530 & 59.061 \\
  & 4 events & 2651 & 1 & 12.172 & 12.172 & 235.825 & 59.217 \\
  & 4 events at different times & 2079 & 1 & 11.153 & 11.153 & 235.355 & 59.146 \\
\midrule
\multirow{3}{*}{Infection (Figure~\ref{fig:epidemic_r0_94})}
  & $R_0=2.6$ & 8840 & 500 & 5826.840 & 11.654 & 1798.114 & 1554.886 \\
  & $R_0=3.07$ & 8830 & 500 & 5630.155 & 11.260 & 1741.606 & 1520.793 \\
  & $R_0=6.0$ & 8818 & 500 & 4949.313 & 9.899 & 1718.252 & 1480.354 \\
\midrule
\multirow{4}{*}{Vaccination (Table~\ref{tbl:vaccine_configuration_94})}
  & VaccA & 16495 & 500 & 8991.828 & 17.984 & 2246.849 & 2057.299 \\
  & VaccB & 20840 & 500 & 11346.181 & 22.692 & 2250.278 & 2318.228 \\
  & VaccC & 23355 & 500 & 16727.571 & 33.455 & 2862.153 & 2512.703 \\
  & VaccD & 12542 & 500 & 7620.026 & 15.240 & 1750.923 & 1883.745 \\
\midrule
\multirow{4}{*}{Vaccination and restriction (Figure~\ref{fig:vaccination_restriction_results_94})}
  & $R_0=2.6$; restriction $=0.50$ & 9011 & 500 & 4008.347 & 8.017 & 1983.816 & 1743.626 \\
  & $R_0=2.6$; restriction $=0.75$ & 8320 & 500 & 4566.876 & 9.134 & 1681.857 & 1472.081 \\
  & VaccC; restriction $=0.50$ & 16722 & 500 & 9279.227 & 18.558 & 2628.414 & 2313.595 \\
  & VaccC; restriction $=0.75$ & 15429 & 500 & 8139.657 & 16.279 & 2540.851 & 2240.514 \\
\bottomrule
\end{tabular}%
}
\end{table}

\subsection{Controlled Benchmarks for the Restructured Results}

The controlled matrix contains 48 configurations: six routine-mobility
baselines, ten Popular Times suppression or Stage~5 adaptation configurations,
eight Stage~6 flood-aware commuting disruptions, thirteen shelter
configurations, and eleven dialysis configurations (D01--D11, corresponding to
K01--K11 in Table~\ref{tbl:dialysis_results}).  Inpatient-care simulations are
excluded.  The two Stage~6 no-flood configurations are included with routine
mobility rather than duplicated in the disruption table.

Every configuration is executed for repetitions 0--4 on the machine described
above and on one repository commit, giving 240 controlled runs.  The runner
schedules up to four independent simulations concurrently.  This five-repetition
rule also applies to configurations whose scenario outcomes are deterministic:
repeating them does not create domain uncertainty, but it does measure runtime
and memory variation under the shared-machine batch workload.  Runtime is
measured with a monotonic high-resolution clock around each simulation loop.
Peak working set is the maximum resident memory reported for each process by
Windows, converted to MiB.  Accordingly, the new timings include CPU, memory,
and storage contention among concurrently executing simulations and should not
be interpreted as isolated single-process runtimes.
Maximum blob count is calculated from the global blob-count series written for
every simulation.  Region- and node-level blob series are deliberately omitted
to avoid very large logging outputs.  TraceMalloc is not collected for the new
results.

The generated tables below are refreshed only from complete controlled runs.
A run is complete only when its metadata, domain validation artifacts, global
blob-count series, runtime, peak working set, and maximum blob count are all
present and mutually consistent.  Missing observations remain explicit rather
than being filled with runtimes from the concurrent production batches.  The
exact commit and individual-run measurements are retained in the accompanying
compact CSV summaries.  After validation and aggregation, the dedicated
performance runs retain only metadata, the global blob-count series, validation
reports, and the small domain tables required by the resume checks.  Redundant
per-step mobility, demand, and OD files are removed because the controlled batch
is used for computational measurements rather than as a replacement for the
separately preserved domain-result batches.

% BEGIN GENERATED PERFORMANCE TABLES
\begingroup
\footnotesize
\setlength{\tabcolsep}{2.6pt}
\begin{table}[htbp]
\centering
\caption{Controlled performance for routine mobility.}
\label{tab:quantitative-performance-routine}
\begin{tabularx}{\textwidth}{@{}>{\raggedright\arraybackslash}Xrrrrrrr@{}}
\toprule
Configuration & Regions & Cycles & Reps. & \begin{tabular}[c]{@{}c@{}}Mean max.\\blobs\end{tabular} & \begin{tabular}[c]{@{}c@{}}Mean runtime\\(s)\end{tabular} & s/cycle & \begin{tabular}[c]{@{}c@{}}Mean peak\\working set\\(MiB)\end{tabular} \\
\midrule
R01 Popular Times only & 13 & 56 & 5 & 1738.0 & 234.736 & 4.1917 & 218.0 \\
R02 Popular Times + probabilistic Lévy & 13 & 56 & 5 & 1832.2 & 286.372 & 5.1138 & 256.9 \\
R03 Popular Times + attendance Lévy V2 & 13 & 56 & 5 & 6962.6 & 454.020 & 8.1075 & 407.2 \\
R04 Popular Times only & 94 & 56 & 5 & 10814.0 & 1893.626 & 33.8147 & 623.2 \\
R05 Popular Times + probabilistic Lévy & 94 & 56 & 5 & 12315.0 & 2916.088 & 52.0730 & 942.9 \\
R06 Popular Times + attendance Lévy V2 & 94 & 56 & 5 & 44580.6 & 3432.490 & 61.2945 & 1824.6 \\
\bottomrule
\end{tabularx}
\end{table}

\begin{table}[htbp]
\centering
\caption{Controlled performance for popular times suppression and stage 5 adaptation.}
\label{tab:quantitative-performance-adaptation}
\begin{tabularx}{\textwidth}{@{}>{\raggedright\arraybackslash}Xrrrrrrr@{}}
\toprule
Configuration & Regions & Cycles & Reps. & \begin{tabular}[c]{@{}c@{}}Mean max.\\blobs\end{tabular} & \begin{tabular}[c]{@{}c@{}}Mean runtime\\(s)\end{tabular} & s/cycle & \begin{tabular}[c]{@{}c@{}}Mean peak\\working set\\(MiB)\end{tabular} \\
\midrule
A01 POI flood; demand suppression & 13 & 56 & 5 & 1832.0 & 277.396 & 4.9535 & 256.0 \\
A02 Home flood; demand suppression & 13 & 56 & 5 & 1870.8 & 297.723 & 5.3165 & 362.7 \\
A03 Home + POI flood; demand suppression & 13 & 56 & 5 & 1837.0 & 278.279 & 4.9693 & 300.7 \\
A04 Home + POI flood; nearest-POI adaptation & 13 & 56 & 5 & 1743.0 & 241.894 & 4.3195 & 273.6 \\
A05 Home + POI flood; Lévy + nearest-POI adaptation & 13 & 56 & 5 & 1838.0 & 294.175 & 5.2531 & 314.8 \\
A06 POI flood; demand suppression & 94 & 56 & 5 & 12316.6 & 2991.751 & 53.4241 & 931.9 \\
A07 Home flood; demand suppression & 94 & 56 & 5 & 12590.6 & 5088.846 & 90.8722 & 2163.3 \\
A08 Home + POI flood; demand suppression & 94 & 56 & 5 & 12338.2 & 3347.209 & 59.7716 & 1347.5 \\
A09 Home + POI flood; nearest-POI adaptation & 94 & 56 & 5 & 10873.0 & 2372.580 & 42.3675 & 1367.4 \\
A10 Home + POI flood; Lévy + nearest-POI adaptation & 94 & 56 & 5 & 12393.0 & 3436.666 & 61.3690 & 1701.1 \\
\bottomrule
\end{tabularx}
\end{table}

\begin{table}[htbp]
\centering
\caption{Controlled performance for stage 6 flood-aware commuting.}
\label{tab:quantitative-performance-disruption}
\begin{tabularx}{\textwidth}{@{}>{\raggedright\arraybackslash}Xrrrrrrr@{}}
\toprule
Configuration & Regions & Cycles & Reps. & \begin{tabular}[c]{@{}c@{}}Mean max.\\blobs\end{tabular} & \begin{tabular}[c]{@{}c@{}}Mean runtime\\(s)\end{tabular} & s/cycle & \begin{tabular}[c]{@{}c@{}}Mean peak\\working set\\(MiB)\end{tabular} \\
\midrule
F01 Flooded activity and commute destinations & 13 & 56 & 5 & 6962.6 & 424.237 & 7.5757 & 399.1 \\
F02 Flooded homes & 13 & 56 & 5 & 8583.4 & 506.052 & 9.0366 & 502.1 \\
F03 Flooded homes and destinations & 13 & 56 & 5 & 8191.2 & 454.694 & 8.1195 & 431.6 \\
F04 All-node flood + nearest-POI adaptation & 13 & 56 & 5 & 8168.6 & 477.270 & 8.5227 & 447.1 \\
F05 Flooded activity and commute destinations & 94 & 56 & 5 & 44579.6 & 3425.326 & 61.1665 & 1757.9 \\
F06 Flooded homes & 94 & 56 & 5 & 51588.6 & 7235.407 & 129.2037 & 2153.8 \\
F07 Flooded homes and destinations & 94 & 56 & 5 & 50274.0 & 4703.032 & 83.9827 & 1687.7 \\
F08 All-node flood + nearest-POI adaptation & 94 & 56 & 5 & 50274.4 & 5188.671 & 92.6548 & 1833.9 \\
\bottomrule
\end{tabularx}
\end{table}

\begin{table}[htbp]
\centering
\caption{Controlled performance for shelter simulations.}
\label{tab:quantitative-performance-shelter}
\begin{tabularx}{\textwidth}{@{}>{\raggedright\arraybackslash}Xrrrrrrr@{}}
\toprule
Configuration & Regions & Cycles & Reps. & \begin{tabular}[c]{@{}c@{}}Mean max.\\blobs\end{tabular} & \begin{tabular}[c]{@{}c@{}}Mean runtime\\(s)\end{tabular} & s/cycle & \begin{tabular}[c]{@{}c@{}}Mean peak\\working set\\(MiB)\end{tabular} \\
\midrule
S01 Reference & 94 & 14 & 5 & 3295.0 & 196.948 & 14.0677 & 510.9 \\
S02 Reference with reallocation & 94 & 14 & 5 & 3353.4 & 206.265 & 14.7332 & 513.7 \\
S03 Capacity 0.5x & 94 & 14 & 5 & 2996.2 & 197.413 & 14.1009 & 511.6 \\
S04 Capacity 2x & 94 & 14 & 5 & 4448.4 & 224.489 & 16.0349 & 517.9 \\
S05 Capacity 5x & 94 & 14 & 5 & 3565.2 & 192.981 & 13.7843 & 513.6 \\
S06 Capacity 15x & 94 & 14 & 5 & 3546.6 & 179.772 & 12.8409 & 512.1 \\
S07 Daily response 0.5\% & 94 & 14 & 5 & 2860.6 & 119.093 & 8.5066 & 497.7 \\
S08 Daily response 2\% & 94 & 14 & 5 & 3662.8 & 468.717 & 33.4798 & 545.3 \\
S09 Daily response 5\% & 94 & 14 & 5 & 3826.4 & 1154.556 & 82.4683 & 640.9 \\
S10 Exposure demand 0.25x & 94 & 14 & 5 & 2171.4 & 76.149 & 5.4392 & 487.6 \\
S11 Exposure demand 0.5x & 94 & 14 & 5 & 2738.4 & 119.896 & 8.5640 & 495.7 \\
S12 Largest-shelter failure & 94 & 14 & 5 & 3317.0 & 203.606 & 14.5433 & 513.7 \\
S13 Top-10\% capacity failures & 94 & 14 & 5 & 3148.8 & 196.992 & 14.0708 & 512.9 \\
\bottomrule
\end{tabularx}
\end{table}

\begin{table}[htbp]
\centering
\caption{Controlled performance for dialysis simulations.}
\label{tab:quantitative-performance-dialysis}
\begin{tabularx}{\textwidth}{@{}>{\raggedright\arraybackslash}Xrrrrrrr@{}}
\toprule
Configuration & Regions & Cycles & Reps. & \begin{tabular}[c]{@{}c@{}}Mean max.\\blobs\end{tabular} & \begin{tabular}[c]{@{}c@{}}Mean runtime\\(s)\end{tabular} & s/cycle & \begin{tabular}[c]{@{}c@{}}Mean peak\\working set\\(MiB)\end{tabular} \\
\midrule
D01 Reduced moderate (K01) & 13 & 10 & 5 & 476.6 & 5.267 & 0.5267 & 459.7 \\
D02 Moinhos persistent outage (K02) & 13 & 10 & 5 & 527.8 & 5.203 & 0.5203 & 457.9 \\
D03 Moinhos rapid recovery (K03) & 13 & 10 & 5 & 477.0 & 5.199 & 0.5199 & 460.3 \\
D04 Historical clinic flood (K04) & 13 & 10 & 5 & 475.8 & 5.355 & 0.5355 & 458.7 \\
D05 Synthetic water-source flood (K05) & 13 & 10 & 5 & 476.2 & 5.058 & 0.5058 & 455.4 \\
D06 Combined reduced disruption (K06) & 13 & 10 & 5 & 476.2 & 4.922 & 0.4922 & 455.5 \\
D07 Reduced above capacity (K07) & 13 & 10 & 5 & 573.2 & 5.810 & 0.5810 & 485.4 \\
D08 One high-capacity failure (K08) & 13 & 10 & 5 & 477.2 & 5.233 & 0.5233 & 461.2 \\
D09 Matched small failures (K09) & 13 & 10 & 5 & 476.2 & 5.107 & 0.5107 & 458.4 \\
D10 Complete-network moderate (K10) & 94 & 10 & 5 & 2763.4 & 16.961 & 1.6961 & 490.2 \\
D11 Complete-network combined (K11) & 94 & 10 & 5 & 2820.2 & 17.792 & 1.7792 & 483.6 \\
\bottomrule
\end{tabularx}
\end{table}

\endgroup
% END GENERATED PERFORMANCE TABLES

\subsection{Interpretation and Limitations}

The legacy and controlled tables should be interpreted separately because the
simulator, logging lifecycle, and measurement protocol changed between them.
Both sets of measurements are hardware-dependent and evaluate the case studies
implemented in LODUS; they are not a head-to-head comparison with other
simulation frameworks.  Five controlled repetitions characterize computational
repeatability on the stated machine, not empirical uncertainty in mobility,
flood exposure, shelter demand, or health-service behavior.
